\documentclass[a4paper,11pt]{article}

% --- Packages ---
\usepackage[utf8]{inputenc}
\usepackage[T1]{fontenc}
\usepackage[a4paper,margin=2.5cm]{geometry}
\usepackage{graphicx}
\usepackage{booktabs}
\usepackage{amsmath,amssymb}
\usepackage{hyperref}
\usepackage{natbib}
\usepackage{caption,subcaption}
\usepackage{float}

\title{PBTA Analysis -- Summary Report}
\author{Alon Luboshitz}

\begin{document}

\maketitle

\section{Introduction}

\section{Data Processing and Methods}

\section{Results}
in tf vs os\_status we see high tf has \big(lower) death rate which is wierd. the km plot show differenlty.
this might be due to the fact alot of samples with high tf (1) has lower death rate (why?) but at high tf ratios (0.7-0.9~) 
there is high death, and the tf=1 is baising the analysis. we should separate tf =1 samples and see the results. 

A question im asking is what cancer group is interesting -  this could be tackled in few ways and its an intreseting
question by it self. First i can ask is which group has strong statistics to work on (number of samples and so on).
Then i see two approaches to tackle this question: unsupervised way - which i conducted here.
I want to find "interseting" groups without prior knoweledge or bais based only on the data.
To do so, i looked for variable - first clinical, then genomic that can shed light on groups/sub groups in the data. 
I followed the demographical analysis in ~cite{} and conducted
\section{Conclusions and Future Directions}

\section{References}

\end{document}
